\section{Introduction}

System Health Management (SHM) has evolved from simple red-line alarms and human-initiated responses to a discipline that often includes sophisticated modeling and automated fault recovery recommendations \cite{Aaseng2001abr}. The main goal of modern SHM has been defined as the preservation of a system's ability to function as intended \cite{Rasmussen2008,Johnson2011a}. While in this paper we apply the term SHM to the operational phase of a system's lifespan, in other contexts the term may also encompass design-time considerations \cite{Johnson2011a}. The actual achievement of a system's operational objectives, on the other hand, is under the purview of the fields of control, planning, and scheduling. In this paper, we will generally refer to all of the processes aimed at accomplishing operational objectives as decision making (DM), while using the more specialized terms where necessary.

Historically, the field of SHM has developed separately from DM. The typical SHM functions (monitoring, fault detection, diagnosis, mitigation, and recovery) were originally handled by human operators--and similarly for DM \cite{Aaseng2001abr,Ogata2010}. As simple automated control was being introduced for DM, automated fault monitors and alarms reduced operator workload on the SHM side. Gradually, more sophisticated control techniques were developed for DM \cite{Ogata2010}, while automated emergency responses started handling some of the off-nominal system health conditions \cite{Rasmussen2008}. Capable automated planning and scheduling tools eventually became available for performing more strategic DM \cite{Ghallab2016}. On the SHM side, automated fault diagnosis was added, in some cases coupled with failure prediction (i.e., \textit{prognostic}) algorithms \cite{Lu1979}, as well as with recovery procedures \cite{Avizienis1976}. Still, the two sides largely remained separated. When the concept of Integrated System Health Management became formalized, most interpretations of \textit{integrated} encompassed some interactions between DM and SHM, although practical implementations of such interactions have been limited \cite{Figueroa2018}.

This paper makes the following claims:

\begin{enumerate}

	\item For actively controlled systems, prognostics is not meaningful; for uncontrolled systems prognostics may only be meaningful under a specific set of conditions;
	\item DM should be unified with SHM for optimality and there is a path for doing so;
	\item Automated emergency response should be done separately from unified DM/SHM, in order to provide performance guarantees and dissimilar redundancy.

\end{enumerate}

For the second claim, we intend to show a path towards DM/SHM unification that builds on the latest developments in state space modeling, planning, and control. While in the past limitations in computing hardware and algorithms would have made unified DM/SHM difficult, we believe that the advances of the recent years make it an attainable goal. We also believe that this unified approach is applicable to a broad spectrum of DM, from traditional deterministic planners to complex algorithms computing long-horizon action policies in the presence of uncertainty. We use the term \emph{unification} to emphasize the idea of DM and SHM being done within the same framework, rather than the two being integrated as separate subsystems exchanging information.

Some initial progress towards DM/SHM unification can be found in the earlier work by \citet{Balaban2013a}, \citet{Balaban2018}, and \citeauthor{Johnson2010} (\citeyear{Johnson2010,Johnson2011a}). System health information was also incorporated into DM by others \cite{Bethke2008,Ure2013,Agha-mohammadi2014}. This paper aims to introduce a systematic view on such integration, discuss its benefits, and illustrate how current SHM concepts map into the proposed approach without a loss of functionality or generality.

The paper first defines the categories of systems that are of interest to this study (Section \ref{sec:systems}) and then overviews the prevailing approach to  SHM (Section \ref{sec:shm}). The first claim (on prognostics) is discussed in Section \ref{sec:prognostics}. Section \ref{sec:hadm} discusses the second claim, concerning DM/SHM integration and its benefits, as well as the rationale for the third claim (that DM/SHM should be separate from automated emergency response). Section \ref{sec:conclusions} concludes.
